% !TeX program = lualatex
% !TeX encoding = utf8
% !TeX spellcheck = uk_UA
% !TeX root = ../main.tex

%% --------------------------------------------------------
%% Chapter file: Motion.tex
%% --------------------------------------------------------

\graphicspath{{\currfilebase/Pictures}}

\chapter{Рух зарядженої частинки в~електромагнітних полях}\label{sec:particle-motion}\hypertarget{particle-motion}{}

Колективний опис не скасовує руху окремих частинок. Щоб зрозуміти магнітне утримання,
спочатку треба побачити найпростішу механіку однієї зарядженої частинки. Вона дає природні
просторовий і~часовий масштаби $r_L$ та $\omega_c^{-1}$, після чого стають зрозумілими
дрейфи, магнітне дзеркало і~межі застосовності МГД.
%% --------------------------------------------------------


%% --------------------------------------------------------
\section{Однорідне магнітне поле}
%% --------------------------------------------------------



У~магнітному полі на частинку з~зарядом $q$ діє сила Лоренца $\vect F=(q/c)\,\vect
v\times\vect B$. Вона перпендикулярна швидкості й~роботи не виконує, тому $v$
зберігається. Розкладемо швидкість $\vect v=\vect v_\parallel+\vect v_\perp$. Складова
$v_\parallel$ вздовж поля не змінюється, а~перпендикулярна складова обертається по колу
з~\emphx{циклотронною частотою} і~\emphx{ларморівським радіусом}
% ~~~~~~~~~~~~~~~~~~~~~~~~~~~~~~~~~~~~~~~~~~~~~~~~~~~~~~~~
\begin{equation}\label{eq:larmor}
	\tcbhighmath{\omega_c=\frac{|q|B}{mc},\qquad r_L=\frac{v_\perp}{\omega_c}=\frac{mc\,v_\perp}{|q|B}}.
\end{equation}
% ~~~~~~~~~~~~~~~~~~~~~~~~~~~~~~~~~~~~~~~~~~~~~~~~~~~~~~~~~~~
Траєкторія~--- гвинтова лінія (рис.~\ref{pic:larmor-motion}). Іони й~електрони обертаються
в~протилежні боки: якщо дивитися з~кінця вектора $\vect B$, позитивні частинки обертаються
за годинниковою стрілкою, негативні~--- проти. Частота обертання не залежить від енергії
частинки. Для практичних оцінок
% ~~~~~~~~~~~~~~~~~~~~~~~~~~~~~~~~~~~~~~~~~~~~~~~~~~~~~~~~
\begin{equation*}
	\omega_{ce}\approx1{,}76\times10^{7}\,B[\text{Гс}]~\mathrm{с}^{-1},\qquad
	\omega_{ci}\approx9{,}58\times10^{3}\,\frac ZA\,B[\text{Гс}]~\mathrm{с}^{-1},
\end{equation*}
% ~~~~~~~~~~~~~~~~~~~~~~~~~~~~~~~~~~~~~~~~~~~~~~~~~~~~~~~~~~~~~~~~
а~ларморівські радіуси при $v_\perp=\sqrt{k_{\mathrm B}T/m}$:
$r_{Le}\approx2{,}4\,\sqrt{T_e[\text{еВ}]}/B[\text{Гс}]$~см,
$r_{Li}\approx10^2\,\sqrt{A\,T_i[\text{еВ}]}/(Z\,B[\text{Гс}])$~см. Через малу масу
електрони мають радіус у~$\sqrt{m_i/m_e}\sim40$ разів менший за іонний при однакових
температурах.

\begin{figure}[h!]\centering
	\localinput{larmor-motion.tikz}
	\caption{Гвинтовий рух зарядженої частинки в~однорідному магнітному полі: уздовж
	поля~--- рівномірний, поперек~--- обертання з~частотою $\omega_c$}
	\label{pic:larmor-motion}
\end{figure}

Плазму називають \emphx{замагніченою}, якщо частинка встигає зробити багато обертів між
зіткненнями, $\omega_c\gg\nu$, і~радіус орбіти малий порівняно з~розміром неоднорідності
поля, $r_L\ll L$. Тоді частинка \enquote{прив'язана} до силової лінії, і~рух поперек поля
пригнічений. Оцінимо перенесення. Вздовж поля частинка блукає з~кроком
$\lambda_{\mathrm{пр}}$ і~$D_\parallel\sim v^2/\nu$; поперек крок дорівнює $r_L$,
а~частота кроків~--- $\nu$, тому $D_\perp\sim r_L^2\nu$. Звідси
% ~~~~~~~~~~~~~~~~~~~~~~~~~~~~~~~~~~~~~~~~~~~~~~~~~~~~~~~~
\begin{equation*}
	\frac{D_\perp}{D_\parallel}\sim\left(\frac\nu{\omega_c}\right)^2\ll1,\qquad D_\perp\propto\frac1{B^2}.
\end{equation*}
% ~~~~~~~~~~~~~~~~~~~~~~~~~~~~~~~~~~~~~~~~~~~~~~~~~~~~~~~~~~~~~~~
Це \emphx{класична дифузія} поперек поля. Вона лежить в~основі магнітного утримання:
сильне поле різко зменшує поперечні втрати.



%% --------------------------------------------------------
\section{Дрейфи в~схрещених полях}
%% --------------------------------------------------------



Нехай, окрім $\vect B$, є однорідне електричне поле $\vect E\perp\vect B$. Шукатимемо рух
у~вигляді гвинтового руху на фоні дрейфу $\vect v_d$ із~сталою швидкістю. Для швидкості
$\vect v_d$ сила має зникати: $\vect E+\vect v_d\times\vect B/c=0$. Помножимо векторно на
$\vect B$ і~скористаємося $(\vect v_d\times\vect B)\times\vect B=-B^2\vect v_{d\perp}$:
% ~~~~~~~~~~~~~~~~~~~~~~~~~~~~~~~~~~~~~~~~~~~~~~~~~~~~~~~~
\begin{equation}\label{eq:ExB}
	\tcbhighmath{\vect v_E=c\,\frac{\vect E\times\vect B}{B^2}}.
\end{equation}
% ~~~~~~~~~~~~~~~~~~~~~~~~~~~~~~~~~~~~~~~~~~~~~~~~~~~~~~~~~~~~
Швидкість дрейфу \emph{не залежить від заряду, маси та енергії частинки}, тому електрони
й~іони дрейфують в~один бік, струму й~розділення зарядів нема. Фізичний зміст: у~системі
відліку, що рухається зі швидкістю $\vect v_E$, електричного поля немає (нерелятивістське
наближення), і~частинка просто гвинтом обертається навколо силової лінії.

Узагальнення: на частинку діє довільна стала сила $\vect F\perp\vect B$ (гравітація,
відцентрова сила, сила від градієнта тиску). Вона викликає дрейф
% ~~~~~~~~~~~~~~~~~~~~~~~~~~~~~~~~~~~~~~~~~~~~~~~~~~~~~~~~
\begin{equation}\label{eq:drift-force}
	\vect v_F=\frac{c}{q}\,\frac{\vect F\times\vect B}{B^2},
\end{equation}
% ~~~~~~~~~~~~~~~~~~~~~~~~~~~~~~~~~~~~~~~~~~~~~~~~~~~~~~~~~~~~~~
який уже \emph{залежить від знака заряду}, і~призводить до струмів та розділення зарядів.
Для $\vect F=q\vect E$ маємо~\eqref{eq:ExB}.

У~неоднорідному полі діють ефективні сили. Усереднена за обертом сила, що виштовхує
частинку в~область слабшого поля, дорівнює $-\mu\nabla B$ (див.~нижче; $\mu$~--- магнітний
момент орбіти), а~при русі вздовж викривленої силової лінії з~радіусом кривизни $R_c$ діє
відцентрова сила $mv_\parallel^2\vect R_c/R_c^2$ (вектор $\vect R_c$ напрямлений від
центра кривизни до частинки). За~\eqref{eq:drift-force} вони породжують
\emphx{градієнтний} та \emphx{відцентровий} дрейфи
% ~~~~~~~~~~~~~~~~~~~~~~~~~~~~~~~~~~~~~~~~~~~~~~~~~~~~~~~~
\begin{equation}\label{eq:gradB-drift}
	\vect v_{\nabla B}=\frac{c\mu}{qB^2}\,\vect B\times\nabla B,\qquad
	\vect v_{R}=\frac{mcv_\parallel^2}{qB^2}\,\frac{\vect R_c\times\vect B}{R_c^2}.
\end{equation}
% ~~~~~~~~~~~~~~~~~~~~~~~~~~~~~~~~~~~~~~~~~~~~~~~~~~~~~~~~~~~~~~~~
За порядком величини $v_d\sim ck_{\mathrm B}T/(|q|BL)=v_T\,r_L/L\ll v_T$. Умова $r_L\ll L$
є умовою застосовності дрейфового наближення, а~знак швидкості залежить від знака заряду.

\begin{figure}[h!]\centering
	\localinput{ExB-drift.tikz}
	\caption{Дрейф зарядженої частинки у~схрещених електричному та магнітному полях}
	\label{pic:ExB-drift}
\end{figure}



%% --------------------------------------------------------
\section{Поляризаційний дрейф}
%% --------------------------------------------------------



Дрейф у~сталому електричному полі є особливим: електрони та іони рухаються в~одному
напрямку. Якщо ж електричне поле змінюється з~часом, інерція частинки створює додатковий
дрейф. Для повільно змінного поля $|\partial_t\vect E|\ll\omega_c E$ маємо
\begin{equation}\label{eq:polarization-drift}
 \tcbhighmath{
 \vect v_{\mathrm{pol}}
 =\frac{c}{\omega_c B}\frac{\partial\vect E_\perp}{\partial t}
 =\frac{mc^2}{qB^2}\frac{\partial\vect E_\perp}{\partial t}.
 }
\end{equation}
На відміну від $\vect E\times\vect B$-дрейфу, цей дрейф залежить від знака заряду і~маси.
Він важливий, коли електричне поле змінюється порівняно повільно з~часом гірообертання,
але вже не є строго сталим. Узагалі дрейфове наближення вимагає $r_L\ll L$ та
$\omega\ll\omega_c$.



%% --------------------------------------------------------
\section{Магнітний момент і~адіабатичний інваріант}
%% --------------------------------------------------------



Обертання частинки по колу~--- це кільцевий струм $I=|q|\omega_c/2\pi$ з~площею контуру
$\pi r_L^2$, тобто магнітний момент (у~СГС $\mu=IS/c$)
% ~~~~~~~~~~~~~~~~~~~~~~~~~~~~~~~~~~~~~~~~~~~~~~~~~~~~~~~~
\begin{equation}\label{eq:mu}
	\mu=\frac{|q|\omega_c}{2\pi c}\,\pi r_L^2=\frac{mv_\perp^2}{2B}.
\end{equation}
% ~~~~~~~~~~~~~~~~~~~~~~~~~~~~~~~~~~~~~~~~~~~~~~~~~~~~~~~~~~~
Момент орбіти напрямлений \emph{проти} поля для частинок обох знаків (вони обертаються
в~протилежні боки, але заряди теж протилежні). Це джерело діамагнетизму плазми
(§\ref{sec:magnetic-confinement}).

Якщо поле змінюється повільно (за час обороту $2\pi/\omega_c$ воно змінюється мало;
в~просторі це означає $r_L\ll L$), величина $\mu$ \emph{зберігається} (адіабатичний
інваріант). Це легко побачити, помітивши, що потік поля крізь орбіту $\Phi=\pi r_L^2B=\pi
m^2c^2v_\perp^2/(q^2B)\propto\mu$. При збільшенні $B$ виникає вихрове електричне поле, яке
прискорює частинку так, щоб потік крізь орбіту залишився сталим (як у~надпровідному
кільці). Інакше кажучи, $\mu$ зберігається з~точністю до $\sim r_L/L$ (строге доведення
випливає зі збереження дії $\oint p\,dq$).

Магнітне поле не виконує роботи, тому повна кінетична енергія зберігається. З~урахуванням
$\mu=\mathrm{const}$
% ~~~~~~~~~~~~~~~~~~~~~~~~~~~~~~~~~~~~~~~~~~~~~~~~~~~~~~~~
\begin{equation}\label{eq:energy-mu}
	\frac{mv_\parallel^2}2+\mu B=\mathrm{const}\quad\Rightarrow\quad m\frac{dv_\parallel}{dt}=-\mu\frac{\partial B}{\partial s},
\end{equation}
% ~~~~~~~~~~~~~~~~~~~~~~~~~~~~~~~~~~~~~~~~~~~~~~~~~~~~~~~~~~~~~~~
де $s$~--- координата вздовж силової лінії. У~області зростання $B$ поздовжня швидкість
зменшується: енергія переходить від поздовжнього руху до поперечного. Результат~--- сила
$-\mu\nabla B$, яка штовхає частинку в~бік слабшого поля.

%% --------------------------------------------------------



%% --------------------------------------------------------


%% --------------------------------------------------------
\section*{Задачі}
%% --------------------------------------------------------

\begin{problem}\label{prb:PlasmaLarmor}
	Електрон з~енергією 100~еВ рухається перпендикулярно полю $B=10^4$~Гс. Знайдіть
	ларморівський радіус і~порівняйте його з~$\lambda_{De}$ для
	$n_e=10^{12}~\mathrm{см}^{-3}$, $k_{\mathrm B}T_e=100$~еВ.
\end{problem}

\begin{problem}\label{prb:PlasmaExB}
	У~схрещених полях $E=1$~В/см і~$B=10^3$~Гс (в~СГС $1~\text{В/см}=1/300$~СГС) знайдіть
	швидкість дрейфу. Чи залежить вона від того, чи рухається електрон, чи протон?
\end{problem}

\begin{problem}\label{prb:PlasmaMirror}
	У~магнітній пастці $B_{\min}=2\times10^4$~Гс, $B_{\max}=8\times10^4$~Гс. Знайдіть
	критичний пітч-кут конуса втрат і~частку частинок ізотропного розподілу, що
	потрапляють у~конуси втрат.
\end{problem}

%% --------------------------------------------------------
\section*{Контрольні питання}
%% --------------------------------------------------------

\begin{Questions}
\begin{quizlist}
\item Які фізичні величини визначають ларморівський радіус і~циклотронну частоту?
		  Чому іони й~електрони обертаються в~протилежні боки?

\item Чому дрейф у~схрещених $\vect E$ і~$\vect B$ не залежить від заряду та маси,
		  а~градієнтний дрейф залежить?

\item Що таке магнітний момент частинки? Чому він є адіабатичним інваріантом? Коли
		  порушується?

\item Як працює магнітне дзеркало? Що таке конус втрат?

\end{quizlist}
\end{Questions}
